\begin{table}[ht]
  \footnotesize
  \caption{Complete paradigms of lexically irregular nouns. Direct-plural variants of \ird{študent} are listed on separate lines.}
  \ifparadigmappendix\else\label{tab:nm-irregular-paradigms}\fi
  \begin{tabularx}{\textwidth}{l l L L}
    \toprule
    {\sc noun (cl.)} & {\sc case} & {\sc sg} & {\sc pl} \\
    \midrule
    \ird{pres} (I) & \textsc{dir} & \ird{pres} & \ird{préše} \\
     & \textsc{trs} & \ird{presa} & \ird{préšov} \\
     & \textsc{ind} & \ird{presam} & \ird{préšac} \\
    \addlinespace
    \ird{tóm} (I) & \textsc{dir} & \ird{tóm} & \ird{tome} \\
     & \textsc{trs} & \ird{toma} & \ird{tomóv} \\
     & \textsc{ind} & \ird{tomem} & \ird{tomac} \\
    \addlinespace
    \ird{študent} (mixed) & \textsc{dir} & \ird{študent} & \shortstack[l]{\ird{študentí}\\\ird{študencí}} \\
     & \textsc{trs} & \ird{študenta} & \ird{študentóv} \\
     & \textsc{ind} & \ird{študencí} & \ird{študencím} \\
    \bottomrule
  \end{tabularx}
\end{table}
